\documentclass[11pt]{article}
\usepackage{weekly}
\begin{document}
\weekheader{2}{06/06 \textendash{} 12/06/2026}{Khảo sát RAG và xây dựng Baseline Simple}{w2col}{Đã hoàn thành}

\section{Mục tiêu tuần}
Tuần thứ hai tập trung vào việc khảo sát có hệ thống các phương pháp RAG, lựa chọn thiết kế
phù hợp với đặc thù miền thủ tục hành chính, và hiện thực hoá một đường cơ sở (baseline) hoàn
chỉnh chạy được trên giao diện: từ câu hỏi của người dùng cho tới câu trả lời có tích hợp
trích dẫn nguồn.

\section{Công việc đã thực hiện}

\paragraph{Khảo sát các phương pháp RAG: }
Nhóm đối chiếu các họ kỹ thuật theo từng khía cạnh của pipeline (truy hồi, viết lại truy vấn,
lọc ngữ cảnh, sinh câu trả lời) để rút ra lựa chọn thiết kế cho miền. Kết quả khảo sát được
tổng hợp ở Bảng~\ref{tab:w2-survey}.

\begin{table}[H]
\centering\small
\renewcommand{\arraystretch}{1.25}
\begin{tabular}{@{}lp{8.4cm}@{}}
\toprule
Khía cạnh & Lựa chọn cho miền thủ tục hành chính \\
\midrule
Truy hồi & Lai ghép (Hybrid): kết hợp tìm kiếm ngữ nghĩa bằng vector và tìm kiếm từ khoá toàn văn. \\
Viết lại truy vấn & Viết lại theo lịch sử hội thoại và giải nghĩa từ viết tắt; bỏ qua truy hồi với câu chào hỏi. \\
Lọc ngữ cảnh & Cổng lọc độ liên quan (Relevance Gating) bằng LLM thay vì lấy cứng Top-$k$. \\
Sinh câu trả lời & Trích dẫn nội dòng kèm tô sáng đoạn bằng chứng; ràng buộc chỉ dùng ngữ cảnh đã truy hồi. \\
\bottomrule
\end{tabular}
\caption{Tổng hợp khảo sát phương pháp RAG và lựa chọn thiết kế cho miền.}
\label{tab:w2-survey}
\end{table}

\paragraph{Thiết kế prompt theo đặc thù miền (Domain Prompting): }
Nhóm xây dựng một prompt hỏi đáp (QA prompt) tiếng Việt với các ràng buộc chặt theo nghiệp vụ
hành chính: chỉ trả lời dựa trên ngữ cảnh đã truy hồi và không bịa thêm từ kiến thức ngoài;
giữ nguyên văn tên giấy tờ kèm mã biểu mẫu (CC01, DC02) và số hiệu văn bản (Thông tư
17/2024/TT-BCA); trình bày thành phần hồ sơ dưới dạng danh sách kèm số lượng bản chính và bản
sao; nêu kèm căn cứ pháp lý khi có. Prompt cũng quy định cách xử lý các trường hợp biên: phân
biệt rõ tin nhắn chào hỏi xã giao với câu hỏi tra cứu, và khi thủ tục chỉ khớp một phần với
câu hỏi (lệch về độ tuổi, đối tượng áp dụng hoặc cấp thực hiện) thì phải nêu rõ điểm chưa khớp
thay vì trình bày như khớp hoàn toàn.

\paragraph{Viết lại truy vấn theo ngữ cảnh (Query Rewriting): }
Câu hỏi từ người dùng thường ngắn và mang tính tiếp nối (follow-up), thiếu ngữ cảnh để truy
hồi chính xác. Nhóm hiện thực một bước viết lại truy vấn dựa trên lịch sử hội thoại để tạo ra
một truy vấn độc lập, trọn vẹn về ngữ nghĩa. Bước này đồng thời giải nghĩa các từ viết tắt phổ
biến trong lĩnh vực (CCCD là căn cước công dân, MST là mã số thuế, BHXH là bảo hiểm xã hội,
GPKD là giấy phép kinh doanh, HKTT là hộ khẩu thường trú). Nếu mô hình nhận diện tin nhắn chỉ
là lời chào xã giao, bước này phát tín hiệu bỏ qua truy hồi để hệ thống phản hồi tự nhiên,
tránh tốn tài nguyên tra cứu.

\screenshotfig{img/week2-simple-answer.png}{Giao diện trả lời ở chế độ Simple: câu hỏi của
người dùng và câu trả lời được tổng hợp kèm danh sách nguồn.}{Hỏi đáp ở chế độ đường cơ sở
Simple trên giao diện.}{fig:w2-answer}

\paragraph{Đường cơ sở Simple (Single-pass Pipeline): }
Nhóm hiện thực luồng suy luận một lượt và tích hợp lên giao diện. Luồng vận hành qua bốn bước
nối tiếp (Hình~\ref{fig:w2-rag}): (1) viết lại truy vấn theo ngữ cảnh; (2) truy hồi lai ghép
song song trên chỉ mục vector và chỉ mục toàn văn rồi hợp nhất kết quả; (3) cổng lọc độ liên
quan, trong đó mỗi đoạn truy hồi được một LLM chấm điểm tương đồng với câu hỏi và loại bỏ các
đoạn dưới ngưỡng; (4) ghép các đoạn còn lại làm ngữ cảnh và sinh câu trả lời. Khi cổng lọc loại
bỏ toàn bộ tài liệu, hệ thống phát một tín hiệu thiếu thông tin tường minh thay vì cố trả lời.

\paragraph{Cơ chế trích dẫn và tô sáng (Citation \& Highlighting): }
Câu trả lời được sinh theo cấu trúc hai khối. Khối thứ nhất liệt kê trích dẫn, gán mỗi nguồn
một số tham chiếu định dạng 【n】 kèm cặp tham số \texttt{START\_PHRASE} và \texttt{END\_PHRASE}
trích nguyên văn đoạn bằng chứng. Khối thứ hai là câu trả lời, trong đó số 【n】 được chèn ngay
sau mỗi dữ kiện. Dựa trên đầu ra này, hệ thống đối chiếu cặp cụm từ với tài liệu gốc để tô
sáng (highlight) chính xác đoạn bằng chứng trên giao diện, đồng thời chuyển các thẻ 【n】 thành
liên kết. Khi điểm liên quan tổng thể của ngữ cảnh ở mức thấp, hệ thống hiển thị cảnh báo để
người dùng chủ động kiểm chứng; ngoài ra một tính năng tuỳ chọn cho phép sinh sơ đồ tư duy
(mind map) tóm tắt thủ tục dưới dạng phân cấp.

\begin{figure}[H]
\centering
\resizebox{\textwidth}{!}{%
\begin{tikzpicture}[node distance=0.6cm and 0.8cm]
  \node[gbox] (q) {Câu hỏi\\+ Lịch sử};
  \node[fbox, right=of q] (rw) {(1) Query\\Rewriting};
  \node[fbox, right=of rw] (ret) {(2) Hybrid\\Retrieval};
  \node[abox, right=of ret] (gate) {(3) Relevance\\Gating};
  \node[fbox, right=of gate] (syn) {(4) Generation\\+ Citation};
  \node[dbox, right=of syn] (ans) {Câu trả lời 【n】\\+ Highlight};
  \draw[farr] (q)--(rw); \draw[farr] (rw)--(ret); \draw[farr] (ret)--(gate);
  \draw[farr] (gate)--(syn); \draw[farr] (syn)--(ans);
  \draw[farr, dashed] (gate.north) to[out=90,in=90] node[flab,above]{Tín hiệu thiếu thông tin} (syn.north);
\end{tikzpicture}}
\caption{Đường cơ sở Simple: suy luận một lượt có trích dẫn.}
\label{fig:w2-rag}
\end{figure}

\screenshotfig{img/week2-citation.png}{Panel nguồn: đoạn bằng chứng được tô sáng và số trích
dẫn 【n】 trong câu trả lời trỏ về đúng nguồn tương ứng.}{Trích dẫn nội dòng và tô sáng đoạn
bằng chứng trong nguồn.}{fig:w2-citation}

\section{Kết quả và sản phẩm bàn giao}
\begin{table}[H]
\centering\small
\begin{tabular}{@{}lp{8.7cm}@{}}
\toprule
Thành phần & Mô tả \\
\midrule
\texttt{rag/prompts.py} & Prompt QA theo miền và prompt viết lại truy vấn (tiếng Việt). \\
Reasoning Simple & Luồng một lượt: viết lại, truy hồi lai ghép, lọc liên quan, tổng hợp. \\
Trích dẫn nội dòng & Định dạng 【n】 với tô sáng \texttt{START}/\texttt{END\_PHRASE}; cảnh báo độ liên quan thấp. \\
Báo cáo khảo sát & Bảng đối chiếu các phương pháp RAG và lý do lựa chọn. \\
\bottomrule
\end{tabular}
\end{table}

\section{Khó khăn và hướng xử lý}
\begin{itemize}
  \item Chiến lược Top-$k$ luôn trả về kết quả kể cả khi tài liệu không liên quan: nhóm bổ
        sung cổng lọc độ liên quan để tạo tín hiệu thiếu thông tin một cách tường minh.
  \item Mô hình đôi khi từ chối nhầm các câu hỏi tưởng nằm ngoài phạm vi: prompt được ràng
        buộc phải sử dụng ngữ cảnh đã chọn, không tự phán đoán câu hỏi là lạc đề.
  \item Sai lệch tên giấy tờ và mã biểu mẫu khi tóm tắt: prompt buộc giữ nguyên văn và liệt
        kê đầy đủ, kèm số bản chính và bản sao.
\end{itemize}

\section{Checklist tuần 2}
\begin{itemize}[label={}]
  \cbDone Khảo sát các phương pháp RAG và chọn thiết kế cho miền.
  \cbDone Prompt QA theo miền và prompt viết lại truy vấn (tiếng Việt).
  \cbDone Đường cơ sở Simple trên giao diện (truy hồi lai ghép và cổng lọc).
  \cbDone Cơ chế trích dẫn nội dòng và tô sáng đoạn bằng chứng.
\end{itemize}

\noindent Hướng tuần kế tiếp: nâng cấp lên Agentic RAG (tác tử ReAct hai giai đoạn), tích hợp
API và giao diện, đồng thời minh bạch hoá quá trình suy luận của agent.

\end{document}
